\begin{center}
	\subsection*{Executive Summary}
\end{center}

\noindent Students who work, commute or care for family must fit their classes around fixed commitments, and building a timetable by hand from thousands of valid combinations is slow. UniPlan lets a student describe their week in one sentence and returns three clash-free timetables with a reason for each.

The system follows an engineering pattern: learned models interpret and rank, and verifiable code decides what is allowed. Rules split the sentence into clauses, a Qwen2.5-7B model fine-tuned with QLoRA on a laptop GPU labels each one, and rules read the days, times and amounts; the student confirms the result or answers a short clarifying question. A deterministic scheduler then finds every valid timetable exactly, so a clash is never recommended. A logistic regression with student-specific weights ranks the candidates, and its per-feature contributions become the reasons.

There was no real student data and no participants were recruited, so every sentence, persona, label and test was written by AI models: Qwen generated most of the training data, and fresh Claude agents wrote and answered the final exams. On the final exams, which were committed after the models were frozen and scored once, the parser read \pRtwoC\% of clauses and \pRtwoS\% of whole sentences correctly (keyword rules: \pKeyC\% and \pKeyS\%), and its clarifying questions covered \cCoveredPct\% of its remaining errors. The ranker agreed with independent answers on \rLR\% of timetable pairs, level with equal weights (\rEq\%) within the exam's uncertainty and well ahead of asking Qwen2.5-14B once per pair (\rQwen\%). End to end, a sentence for two to four units becomes three recommendations in a median of \plMedianRec{}~s. The trained ranker scored higher than equal weights, but the paired interval included zero, so its advantage is unproven.

These results show that the method works on AI-written students, not that it helps real ones. Before wider use we recommend a pilot with real students, accuracy checks for students who write in a second language, and care with demand concentration and the privacy of what students write. The same structure could serve engineering scheduling problems such as production planning and crew rostering.
